\documentclass[10pt]{article}
\input{header.tex}

\begin{document}

\section{Линейные операторы и их матрицы}

\begin{definition}
    Линейным оператором (эндоморфизмом) в линейном пространстве $V$ над полем $K$ называется линейное отображение $\A: V \to V$, т.е.
    \[
    \forall \lambda\in K,\; \forall u,v\in V:\quad \A(u+\lambda v)=\A u+\lambda\A v.
    \]
    Множество всех линейных операторов обозначается $\operatorname{End}(V)$.
\end{definition}

\begin{definition}
    Пусть $e=(e_1,\dots,e_n)$ --- базис $V$. Матрицей оператора $\A$ в базисе $e$ называется матрица $A$, столбцами которой являются координатные столбцы образов базисных векторов:
    \[
    \A e_j = \sum_{i=1}^n A_{ij} e_i,\quad j=1,\dots,n.
    \]
    Тогда для любого $x\in V$ с координатным столбцом $[x]_e$ имеем
    \[
    [\A x]_e = A\,[x]_e.
    \]
\end{definition}

\subsection{Изменение матрицы при замене базиса}

Пусть $e$ и $f$ --- два базиса $V$, связанные матрицей перехода $T$: $f_j = \sum_i T_{ij} e_i$, т.е. $[x]_e = T\,[x]_f$.

\begin{theorem}
    Матрицы $A$ (в базисе $e$) и $A'$ (в базисе $f$) одного и того же оператора $\A$ связаны соотношением
    \[
    A' = T^{-1} A T.
    \]
    Такие матрицы называются \textit{подобными}.
\end{theorem}
\begin{proof}
    Для любого $x\in V$ имеем $[\A x]_e = A [x]_e$. Переходя к базису $f$:
    $[\A x]_f = T^{-1} [\A x]_e = T^{-1} A [x]_e = T^{-1} A T [x]_f$.
    С другой стороны, $[\A x]_f = A' [x]_f$. В силу произвольности $x$ получаем $A' = T^{-1} A T$.
\end{proof}

\begin{corollary}
    Следующие величины оператора $\A$ не зависят от выбора базиса (являются инвариантами):
    \begin{itemize}
        \item ранг $\rg\A = \rank A$;
        \item определитель $\det\A = \det A$;
        \item след $\tr\A = \tr A$;
        \item характеристический многочлен $\chi_\A(t) = \det(A - tE)$.
    \end{itemize}
\end{corollary}
\begin{proof}
    Для определителя: $\det(T^{-1}AT)=\det T^{-1}\det A\det T = \det A$. Для следа: $\tr(T^{-1}AT)=\tr(ATT^{-1})=\tr A$. Характеристический многочлен:
    \[
    \det(T^{-1}AT - tE) = \det(T^{-1}(A-tE)T) = \det(A-tE).
    \]
\end{proof}

\section{Собственные значения и собственные векторы}

\begin{definition}
    Число $\lambda \in K$ называется \textit{собственным значением} оператора $\A$, если существует ненулевой вектор $v\in V$ такой, что
    \[
    \A v = \lambda v.
    \]
    Вектор $v$ называется \textit{собственным вектором}, отвечающим собственному значению $\lambda$.
\end{definition}

Условие $\A v = \lambda v$ равносильно $(\A - \lambda \E)v = 0$, т.е. $v \in \Ker(\A-\lambda\E)$. Нетривиальность ядра означает, что оператор $\A-\lambda\E$ вырожден.

\begin{definition}
    \textit{Характеристическим многочленом} оператора $\A$ называется
    \[
    \chi_\A(t) = \det(\A - t\E).
    \]
    Его корни в поле $K$ (если они существуют) являются собственными значениями $\A$.
\end{definition}

\begin{definition}
    Для собственного значения $\lambda$ множество
    \[
    V_\lambda = \Ker(\A-\lambda\E) = \{v\in V \mid \A v = \lambda v\}
    \]
    называется \textit{собственным подпространством}. Его размерность $\gamma(\lambda)=\dim V_\lambda$ называется \textit{геометрической кратностью} $\lambda$.
\end{definition}

\begin{theorem}[О связи кратностей]
    Для любого собственного значения $\lambda$ выполняется неравенство
    \[
    1 \le \gamma(\lambda) \le \alpha(\lambda),
    \]
    где $\alpha(\lambda)$ --- алгебраическая кратность $\lambda$ (кратность корня $\lambda$ в $\chi_\A(t)$).
\end{theorem}
\begin{proof}
    Пусть $\gamma(\lambda)=k$. Выберем базис $v_1,\dots,v_k$ подпространства $V_\lambda$ и дополним его до базиса всего $V$. В этом базисе матрица $\A$ имеет блочный вид
    \[
    \begin{pmatrix}
    \lambda I_k & B \\
    0 & C
    \end{pmatrix}.
    \]
    Характеристический многочлен $\chi_\A(t)=(\lambda-t)^k \det(C-tI_{n-k})$. Следовательно, $\lambda$ --- корень кратности не меньше $k$.
\end{proof}

\begin{theorem}[Линейная независимость собственных векторов]
    Собственные векторы $v_1,\dots,v_m$, отвечающие попарно различным собственным значениям $\lambda_1,\dots,\lambda_m$, линейно независимы.
\end{theorem}
\begin{proof}
    Индукция по $m$. База $m=1$ очевидна. Пусть утверждение верно для $m-1$. Предположим, что $\sum_{i=1}^m \alpha_i v_i = 0$. Применим оператор $\A$:
    \[
    0 = \A\Bigl(\sum \alpha_i v_i\Bigr) = \sum \alpha_i \lambda_i v_i.
    \]
    Вычтем из этого равенства $\lambda_m$ умноженное на исходное:
    \[
    \sum_{i=1}^{m-1} \alpha_i (\lambda_i - \lambda_m) v_i = 0.
    \]
    По индукционному предположению $v_1,\dots,v_{m-1}$ независимы, поэтому $\alpha_i(\lambda_i-\lambda_m)=0$. Так как $\lambda_i\neq\lambda_m$, то $\alpha_i=0$ для $i=1,\dots,m-1$. Тогда из исходной суммы $\alpha_m v_m=0$, а $v_m\neq0$, значит $\alpha_m=0$.
\end{proof}

\section{Диагонализация линейного оператора}

\begin{definition}
    Оператор $\A$ называется \textit{диагонализируемым} (или \textit{имеет простую структуру}), если существует базис пространства $V$, состоящий из собственных векторов $\A$. В этом базисе матрица оператора диагональна:
    \[
    A = \diag(\lambda_1,\dots,\lambda_n),
    \]
    где $\lambda_i$ --- собственные значения (каждое повторяется столько раз, какова его геометрическая кратность).
\end{definition}

\begin{theorem}[Критерий диагонализируемости]
    Оператор $\A$ диагонализируем тогда и только тогда, когда сумма размерностей всех его собственных подпространств равна $\dim V$, т.е.
    \[
    \sum_{\lambda} \gamma(\lambda) = \dim V,
    \]
    и для каждого собственного значения $\gamma(\lambda)=\alpha(\lambda)$ (геометрическая кратность совпадает с алгебраической).
\end{theorem}
\begin{proof}
    Если существует базис из собственных векторов, то, группируя векторы по собственным значениям, получаем прямую сумму собственных подпространств, откуда следует равенство сумм размерностей. Обратно, если сумма размерностей равна $n$, то, объединяя базисы собственных подпространств, получаем базис $V$, состоящий из собственных векторов.
\end{proof}

\begin{corollary}
    Если характеристический многочлен имеет $n$ попарно различных корней в поле $K$, то оператор диагонализируем.
\end{corollary}

\section{Ортогональные и унитарные матрицы. Положительно определённые операторы}

\subsection{Ортогональные и унитарные матрицы}

\begin{definition}
    Вещественная квадратная матрица $Q$ размера $n\times n$ называется \textit{ортогональной}, если $Q^T Q = I_n$ (или, что то же, $Q Q^T = I_n$). Для комплексных матриц аналогично вводится понятие \textit{унитарной} матрицы: $U^* U = I_n$, где $U^* = \overline{U}^T$ — эрмитово сопряжённая.
\end{definition}

\begin{theorem}[Свойства ортогональных (унитарных) матриц]
    \begin{enumerate}
        \item Ортогональная (унитарная) матрица сохраняет скалярное произведение: $(Qx, Qy) = (x, y)$ для любых векторов $x, y$ (в комплексном случае скалярное произведение эрмитово).
        \item Столбцы ортогональной (унитарной) матрицы образуют ортонормированный базис пространства $\R^n$ ($\mathbb{C}^n$).
        \item Определитель ортогональной (унитарной) матрицы равен по модулю единице: $|\det Q| = 1$.
    \end{enumerate}
\end{theorem}

\begin{proof}
    1. Для вещественного случая: $(Qx, Qy) = (Qx)^T (Qy) = x^T Q^T Q y = x^T y = (x, y)$. Для унитарной: $(Ux, Uy) = (Ux)^* (Uy) = x^* U^* U y = x^* y = (x, y)$.
    
    2. Если $Q = [q_1 \dots q_n]$, то $Q^T Q = I$ означает, что $q_i^T q_j = \delta_{ij}$. Следовательно, $\{q_i\}$ — ортонормированная система из $n$ векторов в $n$-мерном пространстве, т.е. базис.
    
    3. Из $Q^T Q = I$ следует $\det(Q^T Q) = (\det Q)^2 = \det I = 1$, поэтому $|\det Q| = 1$. Для унитарной $U^* U = I$ даёт $|\det U|^2 = 1$.
\end{proof}

\subsection{Положительно определённые операторы}

Пусть $V$ — евклидово (или унитарное) пространство со скалярным произведением $(\cdot, \cdot)$.


\textbf{Определение.}
Оператор $A^*: V \to V$ называется \textit{сопряжённым} к оператору $A$, если для всех $x, y \in V$ выполнено
\[
\langle Ax, y \rangle = \langle x, A^* y \rangle.
\]

\begin{definition}
    Самосопряжённый оператор $\A$ (т.е. $\A^* = \A$) называется \textit{положительно определённым}, если $(\A x, x) > 0$ для любого ненулевого $x \in V$. Если $(\A x, x) \ge 0$ для всех $x$, то $\A$ называется \textit{положительно полуопределённым} (или неотрицательно определённым).
\end{definition}

\begin{remark}
    В конечномерном случае, выбрав ортонормированный базис, мы можем отождествить оператор с его матрицей $A$. Тогда условие самосопряжённости означает $A = A^T$ (симметричность) для вещественного случая и $A = A^*$ (эрмитовость) для комплексного.
\end{remark}


\section{Спектральная теорема для симметричных (эрмитовых) матриц}

\begin{theorem}[Спектральная теорема]
    Пусть $A$ — вещественная симметричная матрица ($A^T = A$) или комплексная эрмитова матрица ($A^* = A$). Тогда:
    \begin{enumerate}
        \item Все собственные значения $A$ вещественны.
        \item Существует ортонормированный базис из собственных векторов $A$.
        \item Матрица $A$ представима в виде $A = U \Lambda U^T$ (в вещественном случае) или $A = U \Lambda U^*$ (в комплексном случае), где $U$ — ортогональная (унитарная) матрица, а $\Lambda = \diag(\lambda_1,\dots,\lambda_n)$ — диагональная матрица собственных значений.
    \end{enumerate}
\end{theorem}

\begin{proof}
    Проведём доказательство для вещественного случая. Комплексный случай аналогичен с заменой транспонирования на эрмитово сопряжение.

    \textbf{Шаг 1. Вещественность собственных значений.}
    Пусть $\lambda$ — собственное значение $A$ с собственным вектором $x$ (возможно, комплексным). Тогда $\lambda \|x\|^2 = \lambda x^* x = x^* (\lambda x) = x^* A x$. Но $x^* A x$ — число, равное своему сопряжённому, так как $(x^* A x)^* = x^* A^* x = x^* A x$ в силу симметричности $A$. Следовательно, $\lambda \|x\|^2$ вещественно, и $\lambda$ вещественно. Более того, можно показать, что собственные векторы можно выбрать вещественными.

    \textbf{Шаг 2. Существование ортонормированного базиса из собственных векторов (индукция по размерности).}
    
    \textit{База:} $n=1$ — очевидно.
    
    \textit{Индукционный переход:} Пусть утверждение верно для всех матриц размера $(n-1)$. Возьмём симметричную матрицу $A$ размера $n\times n$.
    
    У неё есть вещественное собственное значение $\lambda_1$ и соответствующий единичный собственный вектор $u_1$ ($\|u_1\|=1$). Дополним $u_1$ до ортонормированного базиса $\R^n$ векторами $w_2,\dots,w_n$ (например, с помощью процесса Грама–Шмидта). Составим ортогональную матрицу $Q_1 = [u_1 \mid w_2 \dots w_n]$.
    
    Рассмотрим преобразование подобия $Q_1^T A Q_1$. Эта матрица симметрична в силу симметричности $A$: $(Q_1^T A Q_1)^T = Q_1^T A^T (Q_1^T)^T = Q_1^T A Q_1$. Её первый столбец равен $Q_1^T A u_1 = Q_1^T (\lambda_1 u_1) = \lambda_1 e_1$, где $e_1 = (1,0,\dots,0)^T$. В силу симметричности, первая строка также имеет вид $(\lambda_1, 0,\dots,0)$. Следовательно,
    \[
    Q_1^T A Q_1 = \begin{pmatrix}
    \lambda_1 & 0 \\
    0 & B
    \end{pmatrix},
    \]
    где $B$ — симметричная матрица размера $(n-1)\times (n-1)$.
    
    По предположению индукции существует ортогональная матрица $Q_2$ размера $(n-1)$, диагонализующая $B$: $Q_2^T B Q_2 = \diag(\lambda_2,\dots,\lambda_n)$. Тогда итоговая ортогональная матрица
    \[
    Q = Q_1 \begin{pmatrix} 1 & 0 \\ 0 & Q_2 \end{pmatrix}
    \]
    диагонализует $A$: $Q^T A Q = \diag(\lambda_1,\dots,\lambda_n)$. Столбцы $Q$ образуют искомый ортонормированный базис из собственных векторов.
\end{proof}

\begin{corollary}[Спектральное разложение]
    Симметричную матрицу $A$ можно представить в виде суммы проекторов на собственные подпространства:
    \[
    A = \sum_{i=1}^n \lambda_i u_i u_i^T,
    \]
    где $\{u_i\}$ — ортонормированный базис собственных векторов.
\end{corollary}
\begin{proof}
    В разложении $A = Q \Lambda Q^T$ произведение $Q \Lambda Q^T$ можно записать как $\sum_{i=1}^n \lambda_i (Q e_i)(Q e_i)^T = \sum_{i=1}^n \lambda_i u_i u_i^T$, так как $u_i$ — $i$-й столбец $Q$.
\end{proof}

\begin{example}
    Диагонализируем симметричную матрицу $A = \begin{pmatrix} 2 & 1 \\ 1 & 2 \end{pmatrix}$. Её собственные значения: $\lambda_1 = 3$, $\lambda_2 = 1$. Соответствующие ортонормированные собственные векторы: $u_1 = \frac{1}{\sqrt{2}}(1,1)^T$, $u_2 = \frac{1}{\sqrt{2}}(1,-1)^T$. Тогда
    \[
    Q = \frac{1}{\sqrt{2}}\begin{pmatrix} 1 & 1 \\ 1 & -1 \end{pmatrix},\quad Q^T A Q = \begin{pmatrix} 3 & 0 \\ 0 & 1 \end{pmatrix},\quad A = 3u_1u_1^T + 1u_2u_2^T.
    \]
\end{example}


\begin{theorem}[Собственные значения положительно определённых операторов]
    Пусть $\A$ — самосопряжённый оператор в конечномерном евклидовом (унитарном) пространстве. Тогда:
    \begin{enumerate}
        \item Если $\A$ положительно определён, то все его собственные значения строго положительны.
        \item Если $\A$ положительно полуопределён, то все его собственные значения неотрицательны.
    \end{enumerate}
\end{theorem}

\begin{proof}
    По спектральной теореме существует ортонормированный базис из собственных векторов $\A$. Пусть $\lambda$ — собственное значение, а $v$ — соответствующий собственный вектор с единичной нормой. Тогда
    \[
    (\A v, v) = (\lambda v, v) = \lambda (v, v) = \lambda.
    \]
    По определению положительной определённости $(\A v, v) > 0$, следовательно $\lambda > 0$. В полуопределённом случае $(\A v, v) \ge 0$, поэтому $\lambda \ge 0$.
\end{proof}

\begin{corollary}
    Обратно, если все собственные значения самосопряжённого оператора $\A$ положительны, то $\A$ положительно определён. Действительно, разложим любой вектор $x$ по ортонормированному базису собственных векторов $v_i$: $x = \sum_i \alpha_i v_i$. Тогда
    \[
    (\A x, x) = \Bigl(\sum_i \alpha_i \lambda_i v_i, \sum_j \alpha_j v_j\Bigr) = \sum_i |\alpha_i|^2 \lambda_i > 0,
    \]
    если $x \neq 0$.
\end{corollary}

\begin{example}
    Матрица $A^T A$ для любой вещественной матрицы $A$ всегда положительно полуопределена, так как $(A^T A x, x) = (A x, A x) = \|A x\|^2 \ge 0$. Неотрицательность собственных чисел $A$ будет использована в следующей части.
\end{example}

\section{Сингулярные числа и сингулярное разложение (SVD)}

Рассмотрим прямоугольную матрицу $A$ размера $m\times n$ над $\R$ (все обобщается на $\mathbb{C}$).

\begin{definition}
    \textit{Сингулярными числами} матрицы $A$ называются квадратные корни из собственных чисел матрицы $A^T A$ (или $AA^T$), взятые в невозрастающем порядке:
    \[
    \sigma_1 \ge \sigma_2 \ge \dots \ge \sigma_r > 0, \quad \sigma_{r+1}=\dots=\sigma_{\min(m,n)}=0,
    \]
    где $r = \rank A$. (Матрицы $A^T A$ и $AA^T$ симметричны и неотрицательно определены, поэтому их собственные числа неотрицательны.).
\end{definition}

\begin{remark}
    $\rank A^T A = \rank A$, поскольку $A^T A x = 0 \Rightarrow x^T A^T A x = (Ax)^T (Ax) = \|Ax\|^2 = 0 \Rightarrow Ax = 0$, а также $Ax = 0 \Rightarrow A^T A x = 0$, и следовательно размерность пространства решений соответствующих СЛУ совпадают, а значит совпадают и их ранги. Аналогично $\rank A A^T = \rank A^T = \rank A$.
\end{remark}

Связь с собственными числами: если $\lambda_i$ --- ненулевое собственное число $A^T A$, то $\sigma_i = \sqrt{\lambda_i}$.

\begin{theorem}[Сингулярное разложение (SVD)]
    Любую матрицу $A$ размера $m\times n$ ранга $r$ можно представить в виде
    \[
    A = U \Sigma V^T,
    \]
    где:
    \begin{itemize}
        \item $U$ --- ортогональная матрица размера $m\times m$ (столбцы --- левые сингулярные векторы, ортонормированные собственные векторы $AA^T$);
        \item $V$ --- ортогональная матрица размера $n\times n$ (столбцы --- правые сингулярные векторы, ортонормированные собственные векторы $A^TA$);
        \item $\Sigma$ --- прямоугольная диагональная матрица размера $m\times n$ с сингулярными числами на главной диагонали:
        \[
        \Sigma = \begin{pmatrix}
        \sigma_1 & & & \\
        & \ddots & & 0 \\
        & & \sigma_r & \\
        & 0 & & 0
        \end{pmatrix}.
        \]
    \end{itemize}
\end{theorem}

\begin{proof}[Доказательство]
    \textbf{Шаг 1.} Матрица $A^T A$ симметрична и неотрицательно определена. По спектральной теореме для симметричных матриц существует ортонормированный базис $\{v_1,\dots,v_n\}$ пространства $\R^n$, состоящий из собственных векторов $A^TA$. Упорядочим их так, чтобы соответствующие собственные числа $\lambda_i$ убывали:
    \[
    A^TA v_i = \lambda_i v_i, \quad \lambda_1 \ge \lambda_2 \ge \dots \ge \lambda_n \ge 0.
    \]
    Положим $\sigma_i = \sqrt{\lambda_i}$ для $i=1,\dots,n$. Тогда $\sigma_i>0$ для $i\le r$ (где $r$ — ранг $A$) и $\sigma_i=0$ для $i>r$.

    \textbf{Шаг 2.} Для $i=1,\dots,r$ определим векторы
    \[
    u_i = \frac{1}{\sigma_i} A v_i \in \R^m.
    \]
    Покажем, что система $\{u_1,\dots,u_r\}$ ортонормирована:
    \[
    (u_i, u_j) = \frac{1}{\sigma_i\sigma_j} (A v_i)^T (A v_j) = \frac{1}{\sigma_i\sigma_j} v_i^T (A^T A) v_j = \frac{1}{\sigma_i\sigma_j} v_i^T (\lambda_j v_j) = \frac{\lambda_j}{\sigma_i\sigma_j} \delta_{ij} = \frac{\sigma_j^2}{\sigma_i\sigma_j}\delta_{ij} = \delta_{ij}.
    \]
    Таким образом, $\{u_i\}_{i=1}^r$ — ортонормированная система в $\R^m$.

    \textbf{Шаг 3.} Дополним систему $\{u_1,\dots,u_r\}$ до ортонормированного базиса $\R^m$ векторами $u_{r+1},\dots,u_m$ (например, с помощью процесса Грама–Шмидта). Построим матрицы:
    \[
    U = [u_1\ u_2\ \dots\ u_m] \in \R^{m\times m},\qquad V = [v_1\ v_2\ \dots\ v_n] \in \R^{n\times n}.
    \]
    Матрицы $U$ и $V$ ортогональны по построению.

    \textbf{Шаг 4.} Проверим, что $A = U \Sigma V^T$. Для этого вычислим $U^T A V$. Заметим, что
    \[
    (U^T A V)_{ij} = u_i^T A v_j.
    \]
    Если $j > r$, то $\sigma_j = 0$, но это не мешает рассмотрению.
    \begin{itemize}
        \item Для $i \le r$ и $j \le r$:
        \[
        u_i^T A v_j = \frac{1}{\sigma_i} (A v_i)^T A v_j = \frac{1}{\sigma_i} v_i^T A^T A v_j = \frac{1}{\sigma_i} v_i^T (\lambda_j v_j) = \frac{\lambda_j}{\sigma_i} \delta_{ij} = \sigma_i \delta_{ij}.
        \]
        \item Для $i \le r$ и $j > r$: $A^TA v_j = 0$, значит $A v_j = 0$, поэтому $u_i^T A v_j = 0$.
        \item Для $i > r$: вектор $u_i$ ортогонален всем $A v_j$ (так как $A v_j$ лежит в линейной оболочке $\{u_1,\dots,u_r\}$), поэтому $u_i^T A v_j = 0$ для всех $j$.
    \end{itemize}
    Следовательно,
    \[
    U^T A V = \begin{pmatrix}
    \sigma_1 & 0 & \cdots & 0 & 0 & \cdots & 0\\
    0 & \sigma_2 & \cdots & 0 & 0 & \cdots & 0\\
    \vdots & \vdots & \ddots & \vdots & \vdots & & \vdots\\
    0 & 0 & \cdots & \sigma_r & 0 & \cdots & 0\\
    0 & 0 & \cdots & 0 & 0 & \cdots & 0\\
    \vdots & \vdots & & \vdots & \vdots & & \vdots\\
    0 & 0 & \cdots & 0 & 0 & \cdots & 0
    \end{pmatrix} = \Sigma.
    \]
    Умножая слева на $U$ и справа на $V^T$, получаем $A = U \Sigma V^T$.
\end{proof}

\subsection{Свойства сингулярного разложения}

\begin{enumerate}
    \item \textbf{Ранг:} $\rank A = r$ — число ненулевых сингулярных чисел.
    \item \textbf{Норма Фробениуса:} $\|A\|_F = \sqrt{\sum_{i,j} a_{ij}^2} = \sqrt{\sum_{i=1}^r \sigma_i^2}$.
    \item \textbf{Спектральная норма:} $\|A\|_2 = \sigma_1$ (наибольшее сингулярное число).
\end{enumerate}

\subsection{Пример}

Пусть $A = \begin{pmatrix}1 & 1\\ 0 & 1\\ 1 & 0\end{pmatrix}$. Вычислим $A^TA = \begin{pmatrix}2 & 1\\ 1 & 2\end{pmatrix}$. Её собственные числа: $\lambda_1=3$, $\lambda_2=1$, так что сингулярные числа $\sigma_1=\sqrt{3}$, $\sigma_2=1$. Собственные векторы: $v_1=\frac{1}{\sqrt{2}}(1,1)^T$, $v_2=\frac{1}{\sqrt{2}}(1,-1)^T$. Тогда $u_1 = \frac{1}{\sigma_1}A v_1 = \frac{1}{\sqrt{3}}(\sqrt{2}, \frac{1}{\sqrt{2}}, \frac{1}{\sqrt{2}})^T$, $u_2 = \frac{1}{\sigma_2}A v_2 = (\frac{1}{\sqrt{2}}, -\frac{1}{\sqrt{2}}, \frac{1}{\sqrt{2}})^T$. Дополняем до ортонормированного базиса $\R^3$ вектором $u_3$, ортогональным первым двум (например, $u_3 = \frac{1}{\sqrt{6}}(1, -2, 1)^T$ после нормировки). Тогда
\[
U = [u_1\ u_2\ u_3],\quad V = [v_1\ v_2],\quad \Sigma = \begin{pmatrix}\sqrt{3} & 0\\ 0 & 1\\ 0 & 0\end{pmatrix},
\]
и $A = U\Sigma V^T$.

\section{Основные матричные разложения (формулировки)}

В этом разделе приведены формулировки классических матричных разложений без доказательств. Предполагается, что все матрицы вещественные (для комплексных случаев заменяем транспонирование на эрмитово сопряжение).

\subsection{LU-разложение}
\begin{definition}
    \textit{LU-разложением} матрицы $A$ размера $n\times n$ называется её представление в виде произведения нижнетреугольной и верхнетреугольной матриц:
    \[
    A = LU,
    \]
    где $L$ — нижняя треугольная матрица с единицами на диагонали (или, в некоторых вариантах, с ненулевыми диагональными элементами), а $U$ — верхняя треугольная матрица.
\end{definition}

\begin{remark}
    LU-разложение существует (и единственно), если все главные угловые миноры матрицы $A$ отличны от нуля. В общем случае, для невырожденных матриц допустимо LU-разложение после перестановки строк или столбцов ($PAQ = LU$, $P$, $Q$~--- перестановочные матрицы).
\end{remark}

\subsection{Разложение Холецкого}
\begin{definition}
    \textit{Разложением Холецкого} симметричной положительно определённой матрицы $A$ называется её представление в виде
    \[
    A = LL^T,
    \]
    где $L$ — нижняя треугольная матрица с положительными элементами на диагонали. Эквивалентная форма: $A = U^TU$, где $U = L^T$ — верхняя треугольная матрица.
\end{definition}

\begin{remark}
    Разложение Холецкого существует и единственно для любой симметричной положительно определённой матрицы. Оно требует примерно вдвое меньше арифметических операций по сравнению с LU-разложением и является численно устойчивым. Для комплексных эрмитовых матриц используется форма $A = LL^*$, где $L^*$ — эрмитово сопряжённая матрица.
\end{remark}

\subsection{QR-разложение}
\begin{definition}
    \textit{QR-разложением} квадратной невырожденной матрицы $A$ размера $n\times n$ называется её представление в виде произведения ортогональной (унитарной) и верхней треугольной матриц:
    \[
    A = QR,
    \]
    где $Q$ — ортогональная матрица ($Q^T Q = I_n$) в вещественном случае или унитарная ($Q^* Q = I_n$) в комплексном, а $R$ — верхняя треугольная матрица размера $n\times n$.
\end{definition}

\subsection{Полярное разложение}
\begin{definition}
    \textit{Полярным разложением} квадратной матрицы $A$ называется её представление в виде
    \[
    A = HU,
    \]
    где $U$ — ортогональная (унитарныя) матрица, а $H$ — симметричная (эрмитова) матрица.
\end{definition}

\begin{theorem}[Существование]
    Любая квадратная матрица $A$ (вещественная или комплексная) допускает полярное разложение. $H$ определяется единственным образом. Если $A$ обратима, то и матрица $U$ определяется единственным образом.
\end{theorem}

\begin{remark}
    Полярное разложение является аналогом полярной формы комплексного числа $z=re^{i\theta}$, где $H$ представляет собой <<масштабирование>>, а $U$~--- <<поворот>>.
\end{remark}

\end{document}